% ============================================================
%  Poul Martin Møller, "Strøtanker" (Stray Thoughts)
%  Selected aphorisms and short essays, written c. 1830s.
%  First collected in: Efterladte Skrifter, bd. 3,
%  udg. F.C. Olsen, København: C.A. Reitzel, 1843.
%  ENGLISH TRANSLATION — SKELETON
%
%  Translation-only project: a clean modern e-book edition exists
%  (Lindhardt og Ringhof, 2018, ISBN 9788726030204), so NO Danish
%  transcription is produced or published. Translate directly from the
%  e-book (Stroetanker.epub, gitignored). The e-book reprints the 1843
%  Efterladte Skrifter text and preserves the original orthography.
% ============================================================
\documentclass[12pt,a4paper]{article}
\usepackage[utf8]{inputenc}
\usepackage[T1]{fontenc}
\usepackage{libertinus}
\usepackage{libertinust1math}
\usepackage{textalpha}
\usepackage[english]{babel}
\usepackage{enumitem}
\usepackage[protrusion=true,expansion=true]{microtype}
\usepackage{setspace}
\usepackage[margin=1.2in]{geometry}
\usepackage{fancyhdr}
\usepackage{hyperref}
\usepackage{marginnote}

\hypersetup{
  pdftitle={Stray Thoughts (Strøtanker)},
  pdfauthor={Poul Martin Møller},
  hidelinks
}

% Original (2018 ed.) page numbers as marginal notes
\newcommand{\opage}[1]{\marginnote{\footnotesize\textit{[#1]}}}
% Separator between individual stray thoughts
\newcommand{\aphsep}{\par\smallskip\centerline{\textasteriskcentered}\smallskip}

\pagestyle{fancy}
\fancyhf{}
\rhead{Møller, \emph{Stray Thoughts} (1843)}
\cfoot{\thepage}

\setstretch{1.3}

\title{\textbf{Stray Thoughts}\\[4pt]\large (Strøtanker)}
\author{Poul Martin Møller}
\date{Efterladte Skrifter, vol.\ 3, ed.\ F.C.\ Olsen\\
  (Copenhagen: C.A.\ Reitzel, 1843)}

\begin{document}
\maketitle
\thispagestyle{fancy}

\bigskip

\section*{Stray Thoughts on Stray Thoughts}
\label{ch:001}
% Danish: Strøtanker om strøtanker

A flight of thought toward the Eternal is often circumscribed by being written
down on paper.

\aphsep

What I now write, I write solely in order to set my thoughts in motion.
Countless stray thoughts are an infinite introduction to thinking. It is like a
concert that is mere prelude, a poem consisting solely of the invocation of the
Muses and promises of song.

\aphsep

When stray thoughts are set forth in a plain language, many of them, though they
well pleased the one who conceived them, seem to the reader to be nothing. When,
on the other hand, they are clad in splendor, no one says of them that they are
utterly nothing. If they are no substance, they are at least a coat.

\aphsep

Stray thoughts, as fruits of the moment's clearer intuition, are poetic by their
aphoristic form, not scientific. They are the culmination points of thinking ---
a kind of hermaphrodite, half poetry, half prose.

\aphsep

One who writes a multitude of aphoristic reflections may be likened to a
hatmaker who fashions a number of hats of \opage{16}various sizes, of which some
fit one head, others another. Aphorisms have their integration in the author's
subjective circle of thought, and those who occupy themselves with reading them
cannot possibly find them all correct or apt, but some agree with him in one
thing, others in another.

\aphsep

These stray thoughts are really manufactured by a steam engine, since they are
driven by the smoke of Anjer's Liverpool.

\aphsep

The reading of stray thoughts is no labor, but a surrogate for conversation.

\aphsep

The stray thoughts of most authors could, if they stood without selection in the
order in which they had occurred to the author, be divided into three
\textit{partes}: πγευματἱχη, ψυχ, κη, σωματικη, or: Stars, Lilies, and
Cabbage-leaves.

\bigskip

\section*{On Philosophy and Philosophers}
\label{ch:002}
% Danish: Om filosofi og filosoffer

Between leasing and owning, life hovers.

\aphsep

Why can no philosophical system raise itself to that height where its
terminology coincides with the primordial crystallization of language, which is
older than all philosophy? Why must philosophers forever, in their philosophy,
develop or restrict the eternal division of concepts that Nature has made? Is it
because the eternal truth, whose unconscious terminology lies in an unadulterated
language, can never be seen except from one side within a system?

\aphsep

Definitions can sometimes help people to hold fast a concept with precision,
insofar as it can be attached to a couple of concepts that have a fixed meaning
for them; yet such definitions help only where such points of support exist,
just as artificial teeth can be fitted only where there are still some firm ones
to which they may be bound.

\aphsep

In a perfectly, organically developed discourse, every concept is the center of
the whole.

\aphsep

The philosophical free-masters belong to no school and stand in no relation to
any of the philosophical standpoints. They are, as Goethe says, fools on their
own account.

\opage{18}It is a praise for a philosophical work that the Epicureans of the
sciences censure it as dry. The true scientific presentation is only to produce
conviction without arousing emotion.

\aphsep

Complete confusion in the system of cognition can never occur when one is
consistent in presenting his worldview; it occurs only when he is overwhelmed by
reminiscences of an alien worldview. For his system then comes to resemble a
knacker's pit, in which there are nothing but carcasses, and here and there a
half-dead dog crawls about.

\aphsep

Certain truths in philosophy can be grasped only by one whose natural desire has
been broken by fate.

\aphsep

When one has some acquaintance with the great philosophers, one can first
understand the small ones; for the consistency of the latter is not expressed in
their writings, but only intended.

\aphsep

It betrays a great want of a universal overview with respect to the development
of human consciousness to want to disparage literary conversation. It follows
upon a profound system with the same necessity as the hundred echoes of a pistol
shot along the Caucasus mountain range.

\aphsep

As a motto for a philosophical work one may fittingly use Odysseus's words:
τἱ πρῶτον, τἱ δ'ὕστατἱον καταλεξω;

\aphsep

Among his various drives, every individual also has an inclination toward a
certain kind of cognition. Truth can be regarded by the single individual only
from a certain side, and in his individuality already lies the germ of the
philosophy he will form for himself. Yet in \opage{19}the various degrees of his
development his personality will always be traceable. Like a tree, he will draw
spiritual nourishment only from what can be assimilated to his nature. Hence the
disgust with which one who is ever advancing regards the thoughts he himself
wrote down in a bygone time; for he sees, as it were, an embryo, an imperfect and
therefore ugly germ of his present cognition.

\aphsep

Trivial minds think that everything they read is old, because, in order to
understand it, they are obliged to translate it into their trivial mode of
representation. (In a certain sense nothing is new, because everything lies in
the proposition A = A.) Such minds found that Kant's practical reason was =
\textit{conscientia}, that all modern ideas could be translated into Latin. They
believe that ideas are something independent of the given presentation, as though
their insipid manner of apprehension were not just as much a certain
presentation. They stand at so low a level that the most diverse figures flow
into one before their dull sight.

\aphsep

After a rich man, like Kant or Schelling, has cast his great treasure out into
the world, there arise many literary money-changers who exchange the great gold
coins into small change of the same stamp for the people. \textit{Athene} was a
money-changer's shop for the nature-philosophers. ``Kjøbenhavns Skilderie'' deals
out Kant's legacy in small coin, though even these tiny denominations come out
among the people so clipped by Jewish hands that one can scarcely recognize the
stamp.

\aphsep

There are more than one would think who have seen the truth, especially in
isolated clear moments; but they lack the perfected mode of presentation and
therefore possess it only as a subjective opinion.

\opage{20}An essential requisite for philosophical speculation (more essential
than most believe) is a correct concept of the relation between word and thought.

\aphsep

Philosophical productivity cannot coexist with an anxious fear of offending
against usage.

\aphsep

By technical terms one compels the reader to seek the word's meaning in the
system and not to seek the word's meaning in its sound. By the word
\textit{aktuel}, in his Pathology, Sibbern compels the reader to receive a new
concept.

\aphsep

Those who believe that a great philosopher's terminology makes it difficult for
them to enter into his system are often much mistaken. When the presentation is
as it ought to be, comprehension is on the contrary made easier by the fixed
terms. What is difficult for them is really to enter into the new concepts. One
who finds it so hard to learn a philosopher's nomenclature does not complain at
all that, in a play read in a couple of hours, in order to hold fast the
conception of the characters he must learn the persons' names, whose number is
sometimes just as great as the terms one complains of so bitterly. In philosophy
it is not the new sign, but the new articulation of concepts, that demands the
exertion.

\aphsep

Of what use are the inspired, witty, satirical prefaces that Hegel has sent ahead
of his severe writings, forbidding at first acquaintance? --- They render the
same service as the trumpeters who stand blaring before the tents at the
Dyrehave. Without those enticing and much-promising forewords, considerably fewer
would presumably enter into combat with the system itself.

\opage{21}The younger Fichte speaks much of having broken through the barriers of
the Hegelian system. He has broken through them in the manner in which a mouse
breaks through a temple, when it makes itself a hole in the wall and runs out
through it.

\aphsep

If Hegel's terminology is to endure, his successors must apply extraordinary
diligence and force to learning his conceptual designations, roughly in the same
way one learns Linnaeus's botanical nomenclature. But those who have the most
productivity are the least capable of it.

\aphsep

A need to order one's concepts of the understanding leads most people to study
metaphysics. Even Hegel often stirs this need, although he seldom tolerates talk
of a science's usefulness. It can be satisfied without a speculative dialectic in
Hegel's sense.

\aphsep

In the hare, as in many people, lasciviousness is joined with cowardice. ---
Such observations, held fast, might come to constitute one of the pillars in a
pathological system.

\aphsep

The Germans' philosophers, who have so large a reading public, can in their
philosophy furnish the unworked lumps of metal and leave it to others to hammer
them out into small coins that can pass into circulation in the larger circle of
people.

\aphsep

Read old Danish and German writings belonging to various sciences, and feel
wonder at the confidential tone in which they speak with their readers. These
old, honest books form a strange contrast to our century's haughtiness. In olden
days one let intuition and concept \opage{22}inwardly interpenetrate each other;
in our times one holds merely to the concept, because the presentation of
intuition is in part subjective.

\aphsep

The sentimental half-philosophers and half-poets, who like bats move only in
twilight and will not let their movements be illuminated by the sun, become
incensed when one sets assertions against theirs; because the assertions cling
only to their personality, not to the almighty concept.

\aphsep

Small philosophers' conviction that consistency reigns in their presentation
rests on a recollection of the systematic coherence that reigned in the works of
genius on which they lean.

\aphsep

A philosophical system comes to live when it ceases to exist as something
determinate, and comes instead to form a wave in the stream of time.

\aphsep

The demand for originality presupposes a kind of idealistic outlook. One assumes
that philosophy is to be invented. Invent the art of printing!

\aphsep

One who forces upon others a philosophical system for which they feel no need
acts like one who compels a person who is not thirsty to drink.

\aphsep

The opinion that the greatest philosophers greatly disagree with one another
comes in large part from the fact that their readers do not understand how to
read. They then fail to notice that the disharmony is only apparent; they cannot
raise themselves above form and recognize the same concepts in highly different
forms.

\aphsep

\opage{23}A philosophical work in which there are no sentences at all that, torn
out of their context, could be made ridiculous to the common man, can have no
scientific significance.

\aphsep

With a matchless self-confidence Hegel counts on his readers' patience. He is not
in the least troubled by the fear of possibly boring people.

\aphsep

A discourse is popular when at every single point it attaches itself to the
circle of consciousness of the multitude.

\aphsep

Hegel is far from philosophical Epicureanism. He speculates with spiritual
muscularity, without wishing, like Berger and Steffens, to have enjoyment at
every step.

\aphsep

Once I heard a short and pithy philosophical dialogue. A. said: ``There are, upon
my soul, concepts a priori,'' and B. said: ``There are, by my salvation, no
concepts a priori.'' Anyone who knows the real connection of the matter and sees
with what right both assertions can be defended surely does not doubt that these
emphatically pronounced convictions were grounded in the fact that both were
right.

\aphsep

There are those who continue their study of philosophical writings so long, in
order to find the system that can satisfy them, that their capacity for
independent thought more and more disappears. It goes with them as with
Pharaoh's seven lean cows, which ate the seven fat ones and yet did not grow fat.

\aphsep

In the eclectic's philosophy the life of the de-souled systems stirs like maggots
in a carrion-vulture.

\opage{24}It is often poverty of spirit that brings people to diligence. They are
obliged to seek foreign thoughts in order to fill the emptiness in their heads.

\aphsep

Dear Jensen! You are astonished that, in spite of the philosophers' persistent
quarrels, I ascribe reality to philosophy. Do you believe, I will ask you in
reply, that any philosopher could come into dispute with you? You are surely too
modest to expect it. But do you not think, then, that the philosophers' quarrels
must be grounded in a mutual understanding, an agreement; that their quarrels are
facts which presuppose that they have a common ground to stand on?

\aphsep

One must work pettily and in a depressed mood in the service of reflection,
before a total intuition suddenly makes everything clear to one, just as sailors
must warp and tow their ships out through a skerry-belt, so that the brisk breeze
out in the deep may fill their sails.

\aphsep

If the exceedingly strict concept of system were to be made entirely valid, it
would not be possible to understand a systematic discourse; for understanding
could obtain no starting point.

\aphsep

Every giant in philosophy shall fall at the hands of dwarves.

\bigskip

\section*{Science and the Knowledge of Nature}
\label{ch:003}
% Danish: Videnskab og naturerkendelse

The power of imagination, as productive, expresses itself either as a lawless
reproduction of sensuous representations, a false imagination, or as fantasy.

\aphsep

There are three kinds of fullness of bliss: 1) organic well-being, sympathy with
inorganic nature; 2) love of humankind; 3) cognition and art.

\aphsep

People generally believe, as soon as someone talks a great deal, theoretically
well, about a science, that he is a master of it. Fine theories in reviews ---
on the contrary!

\aphsep

To amass dead learning is to adorn oneself so as to appear to have well-filled
limbs.

\aphsep

A trick by which one can give oneself an appearance of learning is to cite
notices without reference. One then imagines, if it is a historical notice, that
the author has read so many books that he cannot keep track of them, or has read
this thing so often that he cannot hit upon the single passage where it stands.
If it is an abstraction, then he can lie himself into an acumen he does not
possess. Some do it involuntarily; but others develop it into a principle.

\aphsep

\opage{26}Linnaeus had no sharp sight and could not see various small infusoria
that are viewed through a microscope, which is also why he makes use of the
curious expression \textit{chaos infusorium}. The Dutch naturalist van Swieten
once wanted to show him some such creatures through a magnifying glass; but
Linnaeus insisted there was nothing to see. At this van Swieten grew so incensed
that he gave Linnaeus a box on the ear. This anecdote I have from Professor
Rathke, who in turn heard it from Wahl, who had it from the mouth of his teacher
Linnaeus himself.

\aphsep

An Epicurean tendency in studies, the need to have satisfaction at once, may
readily be found among critical philologists, whose often rounded-off
investigations yield a self-subsistent result, at least a relatively independent
result.

\aphsep

Men's striving to deny philosophy possibility and reality shows precisely that it
presses itself upon them as essentially necessary. They will not say: ``It
exists, but we do not possess it,'' though with respect to mineralogy or the
Arabic language they show greater resignation and without hesitation say: ``We do
not understand these sciences.''

\aphsep

So long as a man of science hopes at some time to come forward as a writer, he is
tolerant in his judgments of books. When he has entirely given up the thought of
authorship, he becomes mercilessly severe. From his judgments one can tell when
the resignation has set in.

\aphsep

Different tribes of people can resemble one another in foolishness, though the
foolishness has a different appearance. \opage{27}When one delivers a dressing
gown to a Chinese to have its like made, he comes back with a copy that
reproduces even the smallest tear. Europeans often make, in a good translation,
the demand that a meaningless passage shall also be meaningless in the
translation, that nonsense which the author wrote in a fit of absent-mindedness
shall in the translation too continue to be nonsense.

\aphsep

By cultivating the sense for the determination of space one also cultivates the
sense for the determination of time, since time, as an object of representation,
is apprehended under the form of space; a part of time must congeal for the
imagination into a part of space.

\aphsep

Those men of science who feel no great need to order their system of thought must
have either much genius or little thoughtfulness.

\aphsep

That a piece of learning has become cognition is seen from its bursting forth
into other cognitions, just as a plant is seen to have taken root when it puts
forth new shoots.

\aphsep

The merely learned proposition lies motionless in mere memory.

\aphsep

One who reads ontology too early will either not understand it, or apply it
without discernment in his scientific studies. He will apply the most general
determinations, without intermediate terms, to the most empirical concepts.

\aphsep

To call the understanding natural logic is roughly the same as to call sight
natural optics. (By sound logic is meant correct thinking.)

\bigskip

\section*{Existentialism}
\label{ch:004}
% Danish: Eksistentialisme

\begin{center}(and/or dialectical Realism).\end{center}

\aphsep

As a proof of how a concept can be expanded when one one-sidedly pursues it in
its consequences, one could set up a whole moral system in which all good was
derived from the quickness of thinking and all evil from its slowness. For a
certain degree of quickness in thinking is of course required in order to come to
a cognition of the good's worth, and it is, after all, a necessary condition of
morality. Hatred arises from the fact that certain representations become fixed,
cannot be thought through to the end, or cannot finish with their reproductions.
One could also say that the \textit{summum bonum} was the right middle way
between Idealism and Realism.

\aphsep

The anxious elaboration in the enunciation of the concept shows that the concept
is not present in its full force; these are defensive works, evidently thrown up
against attack from without, in mistrust of the concept's own power. Where no
inner life is present, one tries in vain to secure for one's discourse a wretched
existence by means of artificial outworks.

\aphsep

Only when a philosophy has thoroughly penetrated the life of an age, when it is
applied by the so-called popular writers \opage{29}to all the small circumstances
of life, does its whole shape develop; its one-sidedness becomes conspicuous, and
it has completed its course. Something similar happens when the theory of art,
driven in a certain direction, has won itself a sect. The imitators lay bare its
weak sides, they grasp the external, they appropriate the manner --- the one
thing they ought to have thrown away. Their works become an involuntary parody of
the one-sided. The old bird Phoenix must die in order to be born in renewed
shape.

\aphsep

By speculative philosophy one obtains not so much answers to the questions that
common human understanding poses, as one comes to recognize that the questions
themselves, as grounded in concepts that have no truth, must fall away.

\aphsep

A dialectical propaedeutic is a splendid means of formation for the study of the
real sciences and for life's practical activity. Insight into the true essence of
the concept leads to a recognition of the relative validity of opposed modes of
representation and paves the way for forbearance and love in the common life of
human beings.

\aphsep

It would be worth the trouble to investigate whether true philosophy does not
make one disinclined to eloquent discourse. Perhaps the rich, streaming discourse
has its ground in a restless search after the connection of things, which is
satisfied only in fleeting moments and struggles with language to hold fast the
vanishing revelations. Probably deep insight makes one sparing of words; for it
approaches pure knowing, perfect stillness, or death. Plato indeed says that all
philosophizing is a striving \opage{30}after death. These reflections readily
arise in one who has noticed how many profound men speak little, and how
superficial heads can feel themselves pleasantly satisfied by endless prattling
with mouth and pen.

\aphsep

Every human being is allotted in life just as much sorrow and joy as is required
for the awakening of his powers. Fortunate is he whose portion of sorrows is
dealt out in great masses, that they may lift his spirit above the petty. He
whose portion is dealt out in small doses is destined to be small; for only in
great struggle is the hero formed. A life whose greatest sufferings consist in
broken needles is suited to women.

\aphsep

Believe me, the cognition of the coming age lives in many an anticipating mind,
which in its better hours is prompted by its better self to pronounce its
demands; indeed, he may feel it with a shudder, as though his whole present
activity were accommodation and illusion. But he calms himself and forces himself
into agreement with the crowd, since otherwise it seems to him that he would have
to sacrifice all his friends' love and set himself at strife with the whole
world. Yet the strife would not become so great as he fears, if he --- as one who
had died to the world --- developed his deepest and most independent conviction;
then it would surely appear that many kindred strings in others would come to
sound in response. The subjective truth (truthfulness) would then show itself to
lead to the truth.

\aphsep

Those who like to read something mystical into philosophy would, at a lower stage
of formation, have liked to hear ghost stories.

\opage{31}The three periods of human life --- sentimentality, mockery, and sound
feeling --- are a miniature image of the human race's development from
superstition to unbelief, and from unbelief to rational faith.

\aphsep

Certain people believe that the truth ought to be veiled, so that one may not
lose one's faith in human beings. They maintain that one ought to seek to conceal
the truth that great men have had great weaknesses. Fools! Do you know what
discoveries in the realm of the sciences may attach themselves to true
psychological observations? Cicero would perhaps not have been so eloquent had he
not been fickle and pliant of temperament; great actresses would have been less
great had they not been wanton. Should one conceal such things? One need not fear
the illumination of truth; it is enough that one abstains from vivisection.

\aphsep

There are philosophers who cannot get the better of their opponents, because
these stand fast on their earthly standpoint. Could they get them to follow into
the speculative region where they are at home, they would easily vanquish them,
just as Heracles of old vanquished Antaeus when he got his feet lifted from the
earth.

\aphsep

The embittered man is a one-sided idealist who stubbornly regards his subjective
standpoint in society as the absolute.

\bigskip

\section*{Religion, Christianity, and Morality}
\label{ch:005}
% Danish: Religion, kristendom og moral

The wonder of the Christian revelation is so old that, even in those who
historically believe in it, it scarcely awakens more godliness and higher life
than everyday marvels do in most onlookers.

\aphsep

A sermon may have all outward divisions without true coherence, and may lack it,
though it be bound together by inner life.

\aphsep

The words of Scripture, which are meant to give a sermon spirit, are often
degraded to giving it merely body.

\aphsep

Just as a beginning physician finds in himself all diseases, a beginning preacher
finds in himself all faults. Both are in a way right.

\aphsep

Certain narrow minds constantly carry godliness on their lips, without being
either hypocrites or religious. It is really a strong love of earthly goods that
brings them ever to feel with dread that the Eternal One can deprive them of
these treasures. Hence one sees that many peevish, orthodox old men stand ill in
temptation, when their duty \opage{33}is in conflict with their acquisitiveness.
Shakespeare's Henry the Sixth is in this respect masterfully portrayed. In spite
of his anxious religiosity he still cannot bring himself to relinquish the crown
of whose rightfulness he strongly doubts. He is far from a saint, as Schlegel
makes him out to be.

\aphsep

Tears that priests shed in the pulpit do not always stream from a holy source.
The hypocrite bursts into weeping out of moved admiration and enthusiasm over his
own eloquence. These Narcissus-tears spring neither from moral zeal nor from
sorrow, but from a pure interest in the beautiful, from the speaker's platonic
love of himself. Most honest speakers are moved when the theme touches their
personal circumstances, when it, for example, tears open ill-healed wounds that
life gave them. Few are the elect who have so great a heart that it bleeds over
the sins of the world.

\aphsep

Theologians who are not also philosophers readily cast sidelong glances at the
philosophers, and, since they do not find themselves in a position to keep
control over the latter's speculative movements, they live in a constant fear
that some hidden heresy or other may lie in the philosophical systems. Every new
philosophical system has, without exception, been condemned by unphilosophical
theologians. These therefore stand in the same relation to philosophy as hens
that have hatched ducklings, and that now, with anger and anxiety, run about at
the edge of a watering-place where they do not find themselves able to continue
their supervision of their foster-children.

\aphsep

In Christianity's propositions concerning the relation between the world and the
Godhead nothing else can be learned than what a completed \opage{34}speculation
teaches. But indeed something other than what the superficial reflection which
gives itself the name of Enlightenment produces in opposition to Christianity's
mysteries, which stand high above common human understanding.

\aphsep

That guilt immediately draws punishment after it is expressed by the word Ἄἰη. It
is the first immediate cognition of that unity which reflection separates into
guilt and punishment, and which philosophy again makes valid. Also in Paul's
Epistle to the Romans sin is recognized as a punishment for sin.

\aphsep

Invention: A personal God, prepared by speculation, such that not only all the
utterances of the Bible, but every statement in the Catholic, Evangelical, and
Reformed creeds, fits it.

\emph{Note. Can also be used by Mohammedans and Jews.}

\aphsep

Crude orthodox priests must owe the unbelieving authors great thanks; for at
bottom they have no firm conviction, and their faith would little by little fade
away were it not strengthened by their struggle against the enemies of blind
faith.

\aphsep

A devil is unthinkable; for it contradicts reason to assume a devil who rightly
knows what evil is, and yet wills the evil.

\aphsep

A preacher should not despair because the living intuition will not always follow
his pen; it is not to be forced. Let him work confidently with the understanding,
and the intuition comes of itself.

\aphsep

\opage{35}Some, in whom moral concepts have no life, nevertheless use them as
scarecrows in disputes with others.

\aphsep

According to a strict system of identity, the Lord may be regarded as a being who
plays chess with himself. Nature really wages a kind of war with itself in the
various beasts of prey.

\aphsep

``When freedom springs from the power of truth and is therefore dependent on
cognition, how then can the \textit{plebs} become moral, since it cannot raise
itself to cognition?'' --- It cannot raise itself to cognition; but it can be
raised to cognition. Cognition is not the affair of a single man; the men of
science find its forms and highest grounds; but it first comes to its actuality
when it has penetrated all the members of a human community, when, as an
animating principle, it regulates all concrete relations. For this to happen,
speculative cognition must have become immediate in the unscientific multitude.

\aphsep

Sermons that are plainly language-studies are the most offensive of all. Sooner
could one climb up to the moon on an attic ladder than conjure the power of the
Holy Spirit into one's speech by philological care and pedantic turns of phrase.

\aphsep

He does not speak the truth who says that man does not give himself laws; for he
does not know what an I is.

\aphsep

One who has an overstrained morality is young or bad; for his morality differs
from its practice.

\aphsep

Scripture's αληθεια: He who says that he loves God and \opage{36}hates his
brother deceives himself, and there is no truth in him.

\aphsep

A hollow fellow who kisses his old, feeble teacher on the hand and thereby brings
him to burst into tears --- a cruel vivisection.

\aphsep

It is not exactly a lie to pass many different judgments on one and the same
person. Those readily grow angry who do not see that, with all one's love for a
person, one may single out qualities in him which, according to the ordinary
traditional views, they cannot abide.

\aphsep

In an age when it is difficult even for the most upright will to obtain any form
for the supernatural, people snatch at the semblance and delight in deceiving
themselves by the semblance --- imagining, for example, that they believe in
ghost stories. (Small poets' desire to find something extraordinary and
supernatural in the small events of their lives.) On the whole, an unbelieving
age is wont to find significance in the accidental.

\aphsep

That priests do not tolerate deep mockery of life comes from the fact that their
contempt for the earthly is not thorough.

\bigskip

\section*{Aesthetic Investigations}
\label{ch:006}
% Danish: Æstetiske undersøgelser

The goddess of poetry is a great coquette. When one desires her company she does
not come; but when one is least prepared for her visit, one often feels the hot
warmth of her love.

\aphsep

The authors whom the world regards as proud are often the humblest. For their
penitence before the Lord in the private chamber they disdain to let the world
know. The more familiar they become with God, the more they defy the world. They
know that they owe no man an accounting, and the semblance of pride becomes a
barrier that keeps them from amalgamating with the times. Precisely because in
their quiet minds they recognize that they have everything from God, they
therefore declare freely to the world that their deeds are great.

\aphsep

When the popular writer or speaker lacks an expression for a concept, he can
nevertheless, by a cunning turn, drive around the stone.

\aphsep

A number of poets who had begun a work of art and could not afford to give it the
requisite fullness made up for the poverty with a few reflections on life. Such a
work may be likened to a bird in a cabinet of natural curiosities that is
\opage{38}stuffed with hay. Some go about it openly and let such patches pass for
what they are, namely philosophical excursuses. Others have a bad conscience and
seek to give this mechanical connection the semblance of organic necessity by
laying the blame for their thoughts on one of the heroes; but the deception is
nevertheless easily conspicuous.

\aphsep

Whereas the actor (according to the common theory) most easily apprehends the
outward phenomena of a folly foreign to him, the moralist, on the contrary, can
most easily lay bare the faults to which he himself is inclined.

\aphsep

The French speak as they write; but the English write as they speak.

\aphsep

The so-called aesthetic commentators, who believe they can capture the spirit of
antiquity's Latin poets by anatomizing them, contribute thereby little to these
authors' honor. They lay bare precisely the weak side of Horace and Virgil. The
ease of a continuous mechanical analysis of their poetry's rhetorical figures
seems to presuppose a mechanical composition. Only the regularly woven cloth lets
itself be picked apart with a needle and woven together again --- not the plant's
fibers and the body's nerves. The Greek poets cannot so comfortably let
themselves be dismembered.

\aphsep

Genius is not merely a result of natural endowment; it is also begotten by fate.
Hatred gives flashes of genius, just as love does. Therefore the one wronged by
fate looks with a clear gaze and a deep-seeing bitterness into the nature and
essence of life. Shakespeare's bitter profundity can only have been produced by
wrong and injustice.

\aphsep

\opage{39}``King Olaf with the red beard!'' and so forth --- such garish, glaring
colors are used in folk poetry. Kings are painted like a king on a playing card
or woodcut, with a red mustache. King Ahasuerus, etc. Sometimes a king also gets
a blue beard.

\aphsep

A lyric poet enunciates an intuition of his own individuality; from this is seen
the contradiction that lies in a lyric poet's imitating others.

\aphsep

It goes with the good as with the beautiful --- it must be immediately cognized,
and cannot be depicted by colorless concepts.

\aphsep

Just as a weaver can unravel cloth to see how it is fashioned, so a man of science
who feels himself powerfully addressed by a work can carefully examine by what
linguistic means the effect has been produced.

\aphsep

When someone publishes a journal, the literary mob takes it for a poor-box set up
by a highway, into which everyone can put his mite --- indeed, into which naughty
boys can stuff filth and stones.

\aphsep

The poets must, in order to sustain their artist-life, devour one another during
the general famine.

\aphsep

Admiration is necessary for old poets, old harlots, and old dancers.

\aphsep

In poetry the public still stands at the same stage as children, who love the
gaudy woodcuts in bookbinders' windows.

\aphsep

\opage{40}Such a nature as Baggesen's is not so devoid of truth as it appears.
Were the man all falsehood, he would have more inner consistency; but as it is,
the old nature at times breaks through like a lightning-flash.

\aphsep

Self-complacency sometimes expresses itself in this, that during the pauses of
one's speech one lets one's tongue lick the lip a little about its red rim, as if
to taste one's own honey-sweet words.

\aphsep

One of Homer's finest advantages over the more recent poets is the unconscious
presentation of the moving relations of life. The sympathy of blood-kindred is
presented with few, naive utterances, without the heart-gripping being much
emphasized. The deepest feeling is given in an easily advancing narrative,
without dwelling upon it. One does not halt over it to unravel feelings from one
another; but it lies boldly packed together in poetic knots.

\aphsep

The sentimental poetry is a \textit{chaos infusorium}.

\bigskip

\section*{Affectation and Conversation}
\label{ch:007}
% Danish: Affektation og conversation

There is this difference between philosophical conversation and the strict
discourse of science: that the former is intelligible in each single sentence,
whereas in the latter each single part must be understood from the whole. But the
degree of mutual understanding attained by conversation is smaller; one does not
understand the other quite thoroughly. A discourse of that kind may be regarded
as an endeavor to make oneself understood. One seeks to strike the unity of a
certain community's cognition.

The conversationalist therefore seeks more immediate enjoyment and may be vain,
whereas the systematic author can more easily be proud than vain. The
conversationalist must either keep within the limits of the concept that he and
his reading world have in common, or at least proceed from it as a necessary
basis. The systematizer, on the other hand, produces an organic cycle of thought
into which the reader must work his way. --- There is a classical discourse that
combines both.

\aphsep

Conversation is a target-shooting. It depends on whether one hits, or comes near
to, the center of the common circle of consciousness.

\opage{42}The conversationalist cannot be assured that his speech finds an echo in
others until he has tried it by communication, unless he is very conceited. But
when he has become convinced by trial that his subjective mode of expression is
understood by a multitude, he can indeed, with each single utterance, vainly
delight in the effect it will produce. The systematizer, on the other hand, works
toward a total effect.

\aphsep

Reconciliation is often self-deception. ``Now I have made peace with him, the
ass!''

\aphsep

A lie is the poetry that does not come from life. The nearer it comes to life, the
more true; the farther, the more a lie.

\aphsep

Affectation sometimes arises from pride. One who wishes to conceal the effect that
others' speech has upon him is at once untruthful and caught in self-deception.
The latter, because he is not, at the same moment the untruth is uttered,
conscious of it, though by a subsequent reflection he must often be able to
recognize his falseness.

To put on airs.

\bigskip

\section*{Formalism and Stylists}
\label{ch:008}
% Danish: Formalisme og stylistikere

Those who write in the Jewish manner (Heine's, Börne's, Menzel's) seem to be
merry, but their mirth recalls the apparent smile of newborn children, which comes
from their having gripes in the belly.

\aphsep

One ought to make a distinction between a book-maker, who collects and transposes
others' thoughts, and a writer or author.

\aphsep

Clarity in presentation sometimes springs from a lack of the ability to think
coherently.

\aphsep

Novel-readers believe that a man and a woman belong together in Nature's
organization like a matched pair of gloves.

\aphsep

A proof that pleasure in style is a lie is that a boy blushes when one calls his
attention to the fact that he uses ``linger'' for ``wait,'' ``rejoinder'' for
``answer,'' or declaims about silver-white elders.

\aphsep

It is Goethe's most essential merit in style that he is not a stylist. Stylists
are led astray by the circumstance that even great men, like Schiller, could be
stylists. They then choose the easy course of imitating this fault in them.
Socrates says it is fitting for boys --- πλάττειυλόγως.

\bigskip

\section*{The Problem of ``Immortality''}
\label{ch:009}
% Danish: Problemet »Udødelighed«

The will shall not determine what is true.

\aphsep

In a certain sense of the word there are no new proofs of immortality.

\aphsep

Oh, what a mistaken outlook, to assume that immortality does not lie in the
fundamental cognition, that it is at first problematic whether it belongs there or
not.

\aphsep

The narrator some time ago visited an unmarried bookkeeper of his acquaintance,
whom he found in a loud-voiced conversation with a candidate in theology.
Immediately upon his entering, the bookkeeper received him with the following
words: ``It is good you have come; you shall now judge between Ferdinand and me. A
moment ago he fetched a splendid treatise on the immortality of the soul, which he
brings straight from the bookbinder. Now he flatly refuses to lend it to me for a
couple of hours. Is that a friendly way to behave?'' --- The theologian answered:
``You know yourself, Julius! You treat books very badly, and I set great store by
collecting myself a well-conditioned library.'' --- ``You always come out with
that talk, my good Ferdinand, because you once saw me cut a \opage{45}roll of
Dutch tobacco on the binding of a lending-library book. I remember very well, it
was The Murderer in the Mountain Cave, or the Bloody Specter. But today I
expressly promised you that I would treat the book nicely.'' --- ``Oh yes, I'll
vouch for that. You yourself said you would read it while you ate your open
sandwich.'' --- ``There you are right,'' answered the bookkeeper; ``but even if I
had put a little spot on the priceless book, would the damage then have been so
irreparable? I reckon the whole glory has cost you about three marks. But the
reason for that act of friendship --- that you refused me the book --- is,
honestly spoken, your self-love. It would never have vexed me had it not been a
work of this kind that you denied me. I have not thought much about the
immortality of the soul since the time when, as a boy, I had my lesson on it in
Campe's Primer and received, as I still clearly remember, my Excellent-minus for
it. But I am at bottom a religious person; I fully recognize that it is worth the
trouble to get that matter cleared up. And for several years I have resolved to
read, some time when the occasion offered, one good work or another on this
subject. Today I happened to have good leisure, because I am sitting waiting for a
couple of good friends with whom I am to drive at one o'clock to Bellevue to eat
fresh cod. Listen, Julius, lend me the book, then; now I have once got it into my
head to read it today. I am so set on it that I could well take it into my head to
send for it at a bookshop, though you yourself know I am not one of those who throw
money away.'' ``It remains as I have said, I will not do it.''

\aphsep

Their voices relieved one another so quickly that the narrator, who had been
called upon to be their arbitrator, could not manage to make his own voice heard.
At length, however, he got an opportunity to make the remark that it had already
\opage{46}become half past twelve, and that the quarrel was therefore useless,
since the bookkeeper could not possibly finish reading the book in the time that
remained. --- ``Good heavens,'' said the bookkeeper, ``is it so late, and I have
not yet begun my dressing. I would nonetheless most gladly have run through the
famous treatise. I fully perceive the importance of occupying oneself now and then
with the chief truths of religion. So long as the world goes along with one, one
is apt to think little of it; but --- say what one will --- no one can be sure of
himself against tasting such hard blows of fate that it is well he can, with
confidence, lift himself up by the thought of a Beyond. I do not give myself out
for a strong spirit: it would be most dear to me if someone could furnish me a
really palpable proof of the immortality of the soul. If I now do not get another
chance to read a thorough work on it, I have, frankly spoken, nothing really firm
to hold to in dark hours. But whose then is really the fault, my good Ferdinand?
By the way, since we have now once brought this matter up, you, being a
theologian, might just as well set forth for me, quite briefly, the best proofs of
the immortality of the soul, while I strop my razor and take off my beard; but do
not come too near me: I cut myself easily.'' --- ``I find it very difficult,''
said the theologian, ``to impart to you thoughts that were really new to you and
could yet be made clear to you in a moment.'' --- ``There we have the old talk,
Ferdinand! When you will make use of an intelligible language, why should I not
grasp it just as well as you? But we know well enough the clannish spirit that has
possessed you and all your gentlemen colleagues with the Latin and Greek culture.
You want to use your own crow's-tongue in order to keep something \opage{47}for
yourselves. But this caste-business will hardly endure once the more liberal ideas
really break through: Now science is to be popularized; the spirit of the age
demands it. Only take care not to use the technical terms under which the learned
hide their thoughts from laymen, and tell me the plain meaning of it in good broad
Danish. But make haste a little, I am afraid we have the carriage here at any
moment.'' --- Ferdinand answered somewhat embarrassed: ``There are those who
appeal to an intimation, a kind of presentiment of a future higher existence.''
--- ``This point you can skip past, I believe neither in intimations nor in
sympathy; it is precisely a strict, incontrovertible proof that I demand of you.''
--- ``Besides, you doubtless also remember that one is wont to appeal to the
disproportion that often exists between the life-circumstances of the good and of
the bad.'' --- ``Yes, that is well enough: the good must often lack reward for
their self-sacrifice and magnanimity; but who has actually promised them payment
for their deeds? Who tells me that virtue is to be rewarded and vice punished? You
must not think, either, that you can come to me with such loose reasoning as a
village priest may offer his confirmands.'' --- ``One has also roughly the same
proof that we usually in school call the moral, in a more scientific form by
Fichte.'' --- ``Skip the form and report to me the content in the language your
mother taught you; both you and I understand it. But there I see the Vienna
carriage. Now you have here wasted our time repeating old propositions from
Balle's Textbook, instead of furnishing me a stringent proof. Farewell!''

\aphsep

That faith in immortality is an essential moment in cognition shows itself really
clearly only when one resolves entirely \opage{48}to efface this concept from
one's system; then it first becomes evident that it is an indispensable moment in
a consistent and harmonious worldview.

\aphsep

One who champions the immortality of the soul now cannot make his faith
universally valid. The negation must now be carried through, and a new affirmation
must proceed from the denial. We cling fast to the wreck of a sinking ship, others
labor at its destruction. These are two interests.

\aphsep

It leads to mysticism and fanaticism to assume that the soul is not immortal; for
then the denier must strive to assimilate himself with the infinite Indeterminate,
which is the only thing that endures.

\aphsep

The disharmony between the prosaic and the poetic world shall, in the fullness of
time, be reconciled in the true existence. This is a truth that the artists
strongly contest, especially the splendid ones, who feel their minds wholly filled
by their calling. But they do it in a state of lovable naivety, like the savage
who hopes to find again his dear weapons in the realm of the dead.

\aphsep

In a certain sense of the word there are no new proofs of immortality.

\bigskip

\section*{Depth Psychology}
\label{ch:010}
% Danish: Dybdepsykologi

Those wholly dissolute people, whose nature is only mimicry, sometimes have so
rich a nature that they can carry out an assumed character with such force that the
individuals who really have this character cannot carry it through more forcibly.

\aphsep

The incoherent are too poor to be evil.

\aphsep

It is remarkable how much more a person's real, bodily presence makes an
impression on others than his words, read or heard at second hand. Even the
fragmentary utterances of a lightly esteemed companion often work more upon a
thinking person than a well-written book.

\aphsep

In the presentation of psychology there may be found a poetic and a prosaic
atomism. An example of the latter is Treschow's; Garve's individual essays. The
poetic atomism expresses itself in ingenious aphorisms.

\aphsep

Living conviction and a lack of consciousness about the nature and essence of the
categories produce, in conjunction with one another, a high degree of
talkativeness.

\aphsep

\opage{50}A beneficial exposition of psychology should aim at showing the unity of
spiritual life in its manifoldness. Therefore the qualitative differences in it
must always be reduced to quantitative ones.

\aphsep

It lies in the nature of the matter that in letters, somewhat more than in oral
conversation, one must put on one's Sunday clothes.

\aphsep

The actions of the best and of the worst people can be calculated, for they follow
determinate laws.

\aphsep

Man's circle of consciousness consists of the two elements: intuition and concept,
or image and concept. It is once again the same opposition as that between soul
and body. Here again the concept in philosophy is to detach itself, to come
forward with independence.

\aphsep

Does he speak the truth who talks with words he does not himself understand? ---
No! Does he then speak the truth who makes use of words whose concept has been
imparted to him by another, without their standing in any organic connection with
the rest of his cognition? Does he speak the truth who disputes with words he does
not understand, who is governed by his own definitions? --- But this is merely a
play on words; one who believes that he speaks the truth is \emph{bona} fide. ---
No, he deceives himself. He is ashamed to stand naked, in casting away all the
traditional concepts that make up his whole scientific stock. He will not come
forward with his own person, does not believe in his person's infinite depth. If
each single man, without fearing censure for simplicity, judged of things as they
presented themselves to him, then splendid phenomena would come to light. But most
people fear \opage{51}to expose their poverty of spirit. They confirm the confused
in their confusion, by tormenting themselves into their views, from fear of being
regarded as people of scant comprehension. --- Therefore Danish people torment
themselves to speak Germanically; therefore people esteem philosophical phrases
which they do not understand.

\aphsep

One who believes himself insulted by others does not always show love of humankind
by remaining passive. His passivity may come from weakness, and the suppressed
self-feeling may call forth inner embitterment. Also, a long-suppressed
indignation can sometimes get air in an all the more violent explosion.

\aphsep

A lacerated mind must really be ascribed to one who hates himself and abandons
himself to the feeling of his own unworthiness, so that it works destructively upon
his vital force. Instead of improving the I that one condemns, one torments and
destroys it, annihilates it. The enjoyment that lies in this self-torment, this
self-murder, consists in the consciousness that the higher I exercises a punishing
justice.

\bigskip

\section*{Politics and Pedagogy}
\label{ch:011}
% Danish: Politik og pædagogik

One ought to permit the dogmas of the state religion to be contested (also by
officials --- by whom else should they be contested?); for thereby many will
indeed be brought to reject them, but if they are true, they will constantly
return with renewed freedom. One who with freedom returns to their confession is
otherwise convinced of their truth than one who never denied or doubted. When, on
the other hand, dogmas carry with them a kind of outward compulsion, they goad to
polemic. Even if they did not carry such compulsion, they would nonetheless ---
als ein von Wielen stark ausgesprochenes Wort, Widerspruch erregen.

\aphsep

If one did not make use of the Latin language at official examinations, the lack
of dialectic would become more conspicuous.

\aphsep

The usefulness of mathematical study is in particular that the cultivator of
science is accustomed to a certain unpathological thinking (if he needs to be
accustomed to it; for there are those who are by nature all too unfeeling in their
studies).

\aphsep

\opage{53}From young revolutionaries one takes the sting by fawning upon them,
just as the peasant tickles a pig under the chin to make it go stiff.

\aphsep

I do not maintain that everyone, in the course of life, ought to maintain a closed
system of concepts; it might perhaps even give life's development a certain
slowness; but a student ought nonetheless, by way of experiment, to enter into
such a one, and academic instruction ought to aim at that.

\aphsep

There must be loose geniuses who are not definitely subordinated within the body of
the state; for the privileged men of science readily become, by their position,
too inclined to take account of common opinion. It was under other political
conditions that the Reformation could go forth from the university. The privileged
have come into their position by being the spokesmen of a certain system and have
attached themselves firmly to it as the condition of their importance.

\aphsep

One who is to reform the world shall not, by seed, bring forth an entirely new
plant, but, by grafting, improve the old ones.

\aphsep

An old academic teacher once said to me that his lectures would be arranged
otherwise if I were to attend them than in the opposite case: ``for,'' he said,
``I have no wish to make dainties for the young students.'' This, which he naïvely
confessed, came from a mood he has in common with several others; they will not
let their subjectivity shine forth when they speak to their inferiors.

\aphsep

\opage{54}Can a philosophical system become popular? Yes! Philosophy is to be
popularized; but when it has completely become so, then its epoch is over. To
popularize a philosophy is not the work of a single man, but of a whole age; the
age always gradually popularizes the true philosophy, but by transforming its
concepts into opinions. But then it is no longer philosophy; then it is consumed
by \textit{tempus edax rerum}, and the philosophers have made a new advance in
cognition, which the scientific \textit{plebs} has to appropriate.

\aphsep

One can learn firmness of character by giving one's resolutions a kind of
objectivity, in that one, for example, writes them down or communicates them to a
person who becomes a witness to whether they are put into effect or not.

\aphsep

A savage has more human worth than a pedantic scholar who has lost the sense for
all that is higher. The former has respect for prophets and singers. The latter
has entirely rooted out in himself the original cognition of the omnipresence of
the universal life. He rejects every deep glance into the unity of life as
theosophical fanaticism. He persecutes the newer light with hatred, as a member of
the Church of Darkness, the ever-recurring race of the Pharisees.

\aphsep

A vain man may be suited to be a teacher of boys.

\bigskip

\section*{Ahasuerus Stray Thoughts}
\label{ch:012}
% Danish: Ahasverus-Strøtanker

When a flock of boys take turns striking at a barrel on Shrove Monday, it takes
luck to become the one who deals the decisive blow. Thus it takes not only genius
but also luck to become the one who first enunciates the philosophical standpoint.
It must first have been prepared by the endeavors of many predecessors.

\aphsep

Ahasuerus can use this reflection, and add that he has often had it in mind to
strike the last blow at the barrel, but has always put it off.

\aphsep

``Your stupid priests believe that there is an absolute difference between good and
evil, but they do not notice that I stand precisely at the zero point of life's
thermometer.''

\aphsep

Ahasuerus, too, continually has problems for his life, but they are small ones,
such as getting a button sewn on his trousers, loosening his knee-buckle, getting
a new pipe-mouthpiece; nothing but small idyllic joys that he has in common with
the peaceful farmer in a happy state.

\aphsep

The eternal Jew need not pay into the widows' fund. --- How much shall the eternal
Jew pay into the widows' fund, when he is 1800 years older than his wife?

\aphsep

\opage{56}``In a flock of merry young lads I first played a good-natured old fellow
who was young with the young; but when I then saw the mood run through all the
gradations from the jovial tone to the self-satisfied candor and the loving
embraces, on to disharmony --- then my face contracted into an expression of
infinite loathing for the whole of existence, and I exclaimed, with an unspeakable
play of features: Bah! They all dropped their glasses, and it grew still as in a
grave. It was a Medusa's head that I showed them.''

\aphsep

Ahasuerus has such complete consciousness of all his movements and of what is
peculiar in his utterances that nothing is any longer involuntary in him; but from
this it follows that he really always plays a role, for he must, according to plan,
do what ought to be involuntary. Thereby he is put in a position to fool a Jew who
wants to talk after him. He is invited to a party where the Jew first talks after
him, but he then himself talks quite otherwise.

\aphsep

``I, with my detailed knowledge of history, can in any state become whatever I
will; only I must take care not to assert the really historical truth against the
fabricated one.''

\aphsep

``It amuses me to let the human being run on for many years like a machine. I
asserted the one side of a reflection against a priest who was eager to contradict
me. Then he wrote for twenty years against the propounded truth and became a
bishop.''

\aphsep

\textsc{Prognostikon}: There is a kind of young, blond, rustic beauties who more
than others may be compared \opage{57}with the lilies of the field, which today
stand and tomorrow are cast into the oven. They are lovely in their sixteenth year;
but they soon become all too rank in their dimensions and get in their face a
coarse expression that suggests a preponderance of their plant-life. If one wishes
to test whether a lovely blonde belongs to these, then let one lead her into a hot
room, and then all the ethereal and lovely will vanish from her face and her eyes.
This is a kind of trial by fire.

\aphsep

Madam, ``who believe that you, as a calm spectator from your sofa, observe the
doings of the world. Were that true, then you would not have the red glow on your
cheeks.''

\aphsep

Ahasuerus wills nothing. He regards himself as infinitely exalted above those who
will something.

\aphsep

``I can render youthful ardor as well as anyone, if only I do not disturb the
illusion by an involuntary yawn that sometimes comes over me.''

\aphsep

A vigorous fellow who disputes with Ahasuerus arrives at the result that he wills
what he wills merely because he wills it, without being able to state any general
purpose.

\aphsep

His eyes look ``like a windowpane lightly bedewed by a young girl's love-sigh.''

\end{document}
